% === Begin Label Data ===
% ID: 849
% 难度星级: 2
% 标签: 三角函数，追及问题，向量应用
% 备注: 
% 组卷引用次数: 0
% === End  Label Data ===

\begin{problem}{2026}{L}{}{}{解三角形，三角函数}
老虎甲在A地发现野鹿乙在北偏东$15^{\circ}$方向上的B地，立刻以$10\sqrt{3}m/s$的速度进行追捕，与此同时，野鹿乙以$10\sqrt{2}m/s$的速度往北偏东$75^{\circ}$方向逃窜，假设甲、乙都是匀速直线运动，且$AB=500(\sqrt{6}-\sqrt{2})m$，则甲能够一次性捕获乙的最短时间为(\hspace{1cm})
\begin{choices}
\item$60s$
\item$80s$
\item$100s$
\item$120s$
\end{choices}
\end{problem}

\begin{answer}
$C$
\end{answer}

\begin{solutions}
设捕获时间为 $t$ 秒，相遇点为 $C$.

则甲的路程 $AC = 10\sqrt{3}t$,乙的路程 $BC = 10\sqrt{2}t$.

由方位角关系可得 $\angle ABC = 120^{\circ}$.

在 $\triangle ABC$ 中，由正弦定理得
$$
\displaystyle \frac{AC}{\sin \angle ABC} = \displaystyle \frac{BC}{\sin \angle BAC}
$$

代入边长数据并化简得
$$
\displaystyle \frac{10\sqrt{3}t}{\displaystyle \frac{\sqrt{3}}{2}} = \displaystyle \frac{10\sqrt{2}t}{\sin \angle BAC}
$$

解得 $\sin \angle BAC = \displaystyle \frac{\sqrt{2}}{2}$.

因为 $\angle ABC = 120^{\circ}$,所以 $\angle BAC = 45^{\circ}$.

进而求得 $\angle ACB = 180^{\circ} - 120^{\circ} - 45^{\circ} = 15^{\circ}$.

再次利用正弦定理建立方程
$$
\displaystyle \frac{AB}{\sin \angle ACB} = \displaystyle \frac{AC}{\sin \angle ABC}
$$

代入已知条件与三角函数值 $\sin 15^{\circ} = \displaystyle \frac{\sqrt{6}-\sqrt{2}}{4}$ 得
$$
\displaystyle \frac{500\left(\sqrt{6}-\sqrt{2}\right)}{\displaystyle \frac{\sqrt{6}-\sqrt{2}}{4}} = \displaystyle \frac{10\sqrt{3}t}{\displaystyle \frac{\sqrt{3}}{2}}
$$

计算等式两边得
$$
2000 = 20t
$$

解得 $t = 100$.

故甲能够一次性捕获乙的最短时间为 $100s$.

\boxed{C}
\end{solutions}